\uuid{OK0I}
\titre{Pgcd}
\auteur{bodin ; croizat ; sacre}
\organisation{exo7}
\licence{CC BY-NC-SA 4.0}
\datecreate{2022-03-08}
\niveau{L1}
\module{Algèbre}
\chapitre{Arithmétique dans Z}
\sousChapitre{Pgcd, ppcm, algorithme d'Euclide}
\difficulte{2}
\theme{pgcd}
\exo7qcmid{hilisit.arithmetique.008}

\contenu{
\begin{question}
\qclassification{103.03}
    Soit $a>0$ un entier strictement positif dont le reste dans la division euclidienne par $8$ est $r=5$. Quelles sont les affirmations vraies ?
    \begin{answers}
        \bad{$a$ est pair.}
        \good{$a$ est impair.}
        \bad{$a$ est nécessairement divisible par $13$.}
        \good{$(a-5)$ est un multiple de $8$.}
    \end{answers}
    \begin{explanations}
     Puisque le reste dans la division euclidienne de $a$ par $8$ est $5$, on peut écrire $ a = 8k + 5 $, avec $k$ un nombre entier (positif car $a > 0$).\\
     On peut réécrire $a = 8k+5 = 2(4k+2) + 1$ : ainsi $a$ est impair.\\
     Pour $k=0$, on a $a=5$ qui n'est pas divisible par $13$ (ou aussi pour $k=2$ avec $a = 21$ par exemple).\\
     Puisqu'on a $(a-5) = 8k$, cela signifie que $(a-5)$ est bien un multiple de $8$.
    \end{explanations}
\end{question}
}
